\chapter{Desenvolvimento}\label{chap:development}

% TODO: This whole chapter
Neste capítulo, serão descritos os aspectos da implementação prática dos algoritmos descritos no \cref{chap:qmdd} junto da integração com a plataforma de desenvolvimento Ket.

Os algoritmos foram desenvolvidos na linguagem de programação Rust, na versão 1.98.1

% Representação de nodos e arestas como entradas numa tabela para evitar redundância

% Criação de tabela de memoização das operações


\section{Experimentos}\label{sec:experiments}
De forma a validar o desempenho do simulador desenvolvido, diferentes algoritmos foram executados utilizando 3 técnicas de simulação disponíveis na plataforma. Cada execução passou por uma combinação de simulador, algoritmo quântico e número de qubits simulados.

No Ket, estão disponíveis 2 principais simuladores: denso e esparso. Na simulação densa, um circuito de $n$ qubits é simulado utilizando um array de números complexos com $2^n$ posições. Nesse array, cada posição representa a amplitude de um dos $2^n$ possíveis estados do circuito. Assim, cada porta é aplicada individualmente sobre as posições do array que estejam atreladas ao qubit afetado pela porta. Por exemplo, um circuito de 2 qubits no estado \ket{00} é representado pelo array [1, 0, 0, 0].

Já a simulação esparsa se baseia em tabelas que armazenam apenas as amplitudes não nulas do circuito. No Ket, existem duas implementações desse método: uma utilizando tabelas hash e outra utilizando vetores. Ao utilizar tabelas hash, cria-se pares de chave e valor, sendo a chave a string de bits que representa um determinado estado e o valor a sua amplitude. Por exemplo, o estado $\frac{1}{\sqrt{2}}(\ket{00} + \ket{11})$ seria representado pela tabela.

$$
\begin{array}{|c|c|}
  \hline
  \text{Chave} & \text{Valor} \\
  \hline
  00 & \frac{1}{\sqrt{2}} \\
  11 & \frac{1}{\sqrt{2}} \\
  \hline
  % \label{tab:sparse_example}
\end{array}
$$

Em adição a esses 2 simuladores, o desempenho do simulador à base de QMDD recém implementado também foi medido.

Já os algoritmos quânticos simulados foram 3: Grover, FALQON, Estimação de Fase Quântica. [[TODO: Descrever algortimos]]. O objetivo em comparar o desempenho do simulador em diferentes algortimos é identificar em quais circunstâncias simulação por meio de QMDD é mais ou menos eficiente. Isso pois, como dito na TODO: cref, a complexidade de um QMDD varia de acordo com a natureza do circuito simulado.

O Desempenho dos algoritmos foi medido em termos de consumo de tempo e memória. Porém, o consumo de memória não apresentou variação significativa entre os algoritmos e simuladores. [[TODO: maybe leave this part out]].

Para evitar flutuações no tempo de execução, a ordem de cada simulação foi organizada para que os simuladores sejam testados sequencialmente. Ou seja, primeiro foram feitas todas as simulações do QMDD, depois todas do simulador denso e por fim todas do simulador esparso. Assim, questões práticas da implementação de cada simulador em relação à arquitetura do computador utilizado terão menos interferência um no outro. Por exemplo, um simulador que utiliza a memória de forma mais organizada não vai ser afetado por um simulador que deixa a memória do computador fragmentada após sua execução.

Além disso, todas as simulaçãoes de combinações foram realizadas múltiplas vezes, utilizando rodadas de aquecimento para reduzir picos nas rodadas iniciais. O número de execuções foi variado de acordo com a quantia de qubits na simulação: quanto mais qubits, menos rodadas. Isso pois, dado o maior número de qubits, o tempo de execução se torna maior, fazendo com que a variação entre execuções seja menos relevante.

Desse modo, cada combinação de simulador e algoritmo foi executada usando diferentes números de qubits no algoritmo. Para isso, cada combinação iniciou seus experimentos com uma quantia mínima de qubits (1 para Grover, 3 para FALQON e 2 para QPE) e, após cada execução bem sucedida, o número de qubits era incrementado em um. Uma execução falha caso o tempo de execução ultrasse um tempo limite previamente estabelecido.

